\chapter{Wrażliwość współczynników na warunki akwizycji wideo}
\label{ch:wrazliwosc}

Pipeline opisany w~poprzednich rozdziałach produkuje wartości współczynników niezależnie od tego, w~jakich warunkach nagrano wideo. Nie wszystkie współczynniki są jednak równie odporne na te warunki: część zależy od perspektywy kamery, część od częstotliwości klatek, a~niektóre pozostają wiarygodne niezależnie od obu. Rozdział ten bada tę wrażliwość i~pokazuje, że najbardziej zawodne przypadki da się wykryć automatycznie, bez dostępu do pomiarów referencyjnych.

\section{Hipotezy}
\label{sec:hipotezy-p3}

Rozdział weryfikuje trzy hipotezy dotyczące wpływu warunków akwizycji na jakość obliczanych współczynników:

\begin{enumerate}
    \item \textbf{Współczynniki przestrzenne są wrażliwe na perspektywę kamery.} Wartości oparte na geometrii keypointów w~płaszczyźnie obrazu (zwłaszcza wzorzec lądowania stopy) zniekształcają się przy ujęciu innym niż ściśle boczne.
    \item \textbf{Współczynniki temporalne są wrażliwe na częstotliwość klatek (FPS).} Przy niskim FPS kwantyzacja czasu na dyskretne klatki wprowadza istotny błąd względny do metryk czasowych.
    \item \textbf{Niską wiarygodność pomiaru można wykryć automatycznie, bez pomiarów referencyjnych} --- analizując, czy wyliczona wartość mieści się w~zakresie fizjologicznie możliwym.
\end{enumerate}

Każdą hipotezę badam na studium przypadku opartym na konkretnych nagraniach ze~zbioru.

\section{Studium 1: wzorzec lądowania stopy a~perspektywa kamery}
\label{sec:case1-foot-strike}

Wzorzec lądowania stopy (sekcja~\ref{sec:obliczanie-wspolczynnikow}) klasyfikowany jest na podstawie kąta wektora od~pięty do~czubka stopy w~klatce wejścia w~fazę kontaktu. Aby ocenić jego wiarygodność, przeprowadziłem walidację wizualną na trzech biegaczach o~różnych warunkach nagrania: Adamie (film~12, zbiór treningowy), biegaczu z~filmu~10 (zbiór testowy) oraz Janku (nagranie spoza zbiorów modelu). Ponieważ jest to walidacja wizualna samej geometrii kąta stopy, a~nie ocena generalizacji klasyfikatora, użycie filmu ze~zbioru treningowego (Adam) nie wpływa na wnioski. Dla każdego biegacza wyrenderowałem szkielet MediaPipe na klatkach wejścia w~fazę STANCE (po trzy klatki na nogę, rozłożone równomiernie) i~porównałem obliczony kąt z~rzeczywistym ułożeniem stopy widocznym na obrazie. Tabela~\ref{tab:foot-strike-validation} zestawia wyniki.

\begin{table}[htbp]
\centering
\caption{Walidacja wzorca lądowania stopy dla trzech warunków nagrania}
\label{tab:foot-strike-validation}
\small
\begin{tabular}{p{2.2cm}p{4.2cm}p{2.6cm}p{3.6cm}}
\hline
\textbf{Biegacz} & \textbf{Ujęcie} & \textbf{Średni kąt L / P} & \textbf{Werdykt} \\
\hline
Janek    & boczne, kamera prostopadle do toru & $-4^\circ$ / $-3^\circ$ & midfoot, wiarygodne \\
Adam     & boczne, kamera lekko od dołu & $-33^\circ$ / $-12^\circ$ & forefoot; lewa noga zawyżona przez perspektywę \\
Film~10  & pionowe (portrait) & $-97^\circ$ / $-99^\circ$ & bezwartościowe ($|\text{kąt}| > 90^\circ$) \\
\hline
\end{tabular}
\end{table}

Wyniki układają się w~czytelny wzorzec. Przy standardowym ujęciu bocznym (Janek) obliczone kąty zgadzają się z~obserwacją wizualną --- stopa ląduje płasko, a~klasyfikacja midfoot jest wiarygodna. Przy ujęciu bocznym, ale z~kamerą ustawioną lekko od dołu (Adam), prawa noga (bliżej kamery) daje sensowny wynik, natomiast lewa (dalej, silniej dotknięta perspektywą) ma kąt zawyżony --- wektor pięta--czubek jest skrócony w~rzucie, co matematycznie wzmacnia składową pionową. Asymetria $-33^\circ$ wobec $-12^\circ$ jest więc artefaktem perspektywy, nie biomechaniki. Przy ujęciu pionowym (Film~10) wszystkie kąty przekraczają $90^\circ$ w~wartości bezwzględnej, co jest geometrycznie niemożliwe dla biegacza zwróconego w~kierunku biegu --- predykcja jest bezwartościowa.

Studium to potwierdza pierwszą hipotezę: wzorzec lądowania stopy jest wiarygodny przy ujęciu ściśle bocznym i~zawodzi przy ujęciu pionowym lub ukośnym. Warto zauważyć, że ta sama pionowa orientacja Filmu~10, która psuje kąt stopy, obniża również dokładność klasyfikatora faz --- problem ten zaadresowałem korektą proporcji obrazu, której wpływ przeanalizowałem osobno w~sekcji~\ref{sec:ablation-aspect}.

\section{Studium 2: częstotliwość klatek a~precyzja metryk temporalnych}
\label{sec:case2-fps}

Współczynniki temporalne wyliczam z~liczby klatek przypadających na daną fazę (sekcja~\ref{sec:segmentacja}). Każda faza trwa całkowitą liczbę klatek, więc jej czas jest skwantowany do wielokrotności $1/\text{FPS}$. Przy wysokim FPS kwantyzacja jest drobna, ale przy niskim staje się istotnym źródłem błędu.

Rozważmy typowy czas kontaktu z~podłożem GCT~$=250$~ms. Przy 30~FPS jedna klatka odpowiada $\approx$~33~ms, więc kontakt obejmuje około 7,5~klatki, a~błąd $\pm 1$~klatki to $\pm 13\%$. Przy 13,33~FPS (jak w~slow-motion Filmie~7) jedna klatka to $\approx$~75~ms, kontakt obejmuje zaledwie $\approx$~3,3~klatki, a~ten sam błąd $\pm 1$~klatki rośnie do $\pm 30\%$. Im krótsza faza i~niższy FPS, tym większy względny błąd kwantyzacji.

Wniosek ma istotną konsekwencję metodologiczną, spójną z~założeniami przyjętymi przy budowie zbioru (sekcja~\ref{sec:dataset}): nagrania o~niskim FPS są pełnoprawnym materiałem \textit{treningowym} dla klasyfikatora faz (model uczy się pozycji ciała, nie czasu), ale nie nadają się do precyzyjnego wyliczania metryk \textit{temporalnych} na etapie inferencji. FPS należy więc traktować jako parametr jakości nagrania przy obliczaniu współczynników czasowych, a~nie przy treningu.

\section{Synteza: wrażliwość metryk i~możliwość auto-detekcji}
\label{sec:synteza-wrazliwosc}

Tabela~\ref{tab:wrazliwosc-synteza} zestawia wszystkie współczynniki pod kątem wrażliwości na orientację kamery i~na FPS oraz tego, czy daną wadliwość da się wykryć automatycznie.

\begin{table}[htbp]
\centering
\caption{Wrażliwość współczynników na warunki akwizycji i~możliwość automatycznej detekcji}
\label{tab:wrazliwosc-synteza}
\small
\begin{tabular}{p{3.4cm}p{2.8cm}p{2.4cm}p{3.4cm}}
\hline
\textbf{Współczynnik} & \textbf{Orientacja} & \textbf{FPS} & \textbf{Auto-detekcja} \\
\hline
Kadencja            & nie & tak & tak (z~metadanych FPS) \\
GCT                 & nie & tak & tak \\
Czas lotu           & nie & tak & tak \\
Stride length       & nie & tak & tak \\
Duty factor         & nie & tak & tak \\
Kąty stawów         & minimalnie & nie & nie (cichy artefakt) \\
Oscylacja pionowa   & tak (out-of-plane) & nie & częściowo \\
Foot strike         & tak (krytycznie) & nie & tak (próg $45^\circ$) \\
Symetria L/P        & tak (perspektywa) & nie & częściowo (SI~$>$~50\% jako przesłanka) \\
\hline
\end{tabular}
\end{table}

Z~tabeli wynika wyraźny podział. Metryki temporalne są odporne na orientację, lecz wrażliwe na FPS --- a~FPS jest znany z~metadanych pliku, więc problem da się wykryć w~pełni automatycznie. Metryki przestrzenne są odwrotnie: niewrażliwe na FPS, ale wrażliwe na perspektywę. Tu auto-detekcja jest trudniejsza, bo nie ma metadanych mówiących o~poprawności ujęcia --- trzeba wnioskować z~samych wartości. Najtrudniejszy przypadek to kąty stawów, których zniekształcenie perspektywiczne jest „ciche": wartość pozostaje w~pozornie sensownym zakresie, więc nie sposób jej odrzucić bez pomiaru referencyjnego.

\section{Automatyczna detekcja niewiarygodnych pomiarów}
\label{sec:auto-detekcja}

Trzecia hipoteza --- możliwość automatycznego wykrycia niewiarygodnych pomiarów bez ground truth --- znajduje potwierdzenie w~przypadkach, gdzie wartość metryki wykracza poza zakres fizjologicznie możliwy. Zaimplementowałem dwa takie mechanizmy.

\textbf{Flaga \texttt{low\_confidence} dla wzorca lądowania stopy.} W~funkcji obliczającej wzorzec lądowania oznaczam wynik jako niewiarygodny, gdy średni kąt stopy przekracza $45^\circ$ w~wartości bezwzględnej. Próg ten skalibrowałem na podstawie walidacji wizualnej ze~Studium~1: Janek (poprawne ujęcie) osiągał $|\text{kąt}| \le 12^\circ$, Adam (ujęcie graniczne) $12$--$33^\circ$, a~Film~10 (ujęcie pionowe) $82$--$160^\circ$. Próg $45^\circ$ oddziela „poprawne z~marginesem" od „ewidentny artefakt". Gdy flaga się zapala, silnik reguł zwraca ostrzeżenie o~niskiej wiarygodności i~wstrzymuje reguły interpretujące wzorzec lądowania --- nie ma sensu doradzać na podstawie wartości, która jest artefaktem perspektywy.

\textbf{Reguła łączona jakości predykcji.} Drugi mechanizm, opisany w~sekcji~\ref{sec:rules-combined}, wychwytuje sytuacje, w~których wysoki wskaźnik symetrii łączy się z~obniżoną pewnością predykcji --- sygnał, że asymetria L/P jest raczej artefaktem klasyfikatora niż cechą biegacza.

Oba mechanizmy działają bez pomiarów referencyjnych: opierają się wyłącznie na tym, czy wartość mieści się w~zakresie możliwym fizycznie i~czy różne wskaźniki jakości są ze~sobą spójne. To potwierdza trzecią hipotezę w~zakresie metryk, których zniekształcenie wyprowadza wartość poza zakres fizjologiczny. Ograniczeniem pozostają „ciche" artefakty kątów stawów, dla których taka detekcja nie jest możliwa.

\section{Walidacja na nagraniu własnym z~pomiarem referencyjnym}
\label{sec:walidacja-garmin}

Dotychczasowe oceny opierały się na porównaniu wyliczonych wartości z~zakresami z~literatury, ponieważ zbiór nie zawierał nagrań z~równoległym, niezależnym pomiarem. Aby uzupełnić ten brak, zarejestrowałem dodatkowe nagranie własnego biegu na bieżni mechanicznej (oznaczone w~danych jako \textit{Wojtek}, nagranie spoza zbiorów treningowego, walidacyjnego i~testowego), rejestrując jednocześnie dane z~zegarka sportowego Garmin. Zegarek zapisuje sekundowe wartości dynamiki biegu --- kadencję, czas kontaktu z~podłożem, długość kroku --- które posłużyły jako pomiar referencyjny (ground truth) dla wyników pipeline'u. Jest to jedyny w~całej pracy punkt odniesienia pochodzący z~zewnętrznego urządzenia, a~nie z~zakresów literaturowych.

Nagranie zarejestrowałem w~rozdzielczości $1920 \times 1080$ (poziomej) przy 29{,}97~FPS. Pełna aktywność obejmowała rozbieganie i~próby nagrania; do analizy wykorzystałem fragment biegu właściwego ($\approx$~35~s, kadencja ustabilizowana). Porównania dokonałem względem odpowiadającego mu okna danych Garmina --- średnia z~całej aktywności byłaby zaniżona przez wolne rozbieganie i~przerwy. Tabela~\ref{tab:garmin-validation} zestawia wyniki.

\begin{table}[htbp]
\centering
\caption{Porównanie pipeline'u z~pomiarem referencyjnym Garmin (okno biegu właściwego)}
\label{tab:garmin-validation}
\small
\begin{tabular}{lccc}
\hline
\textbf{Metryka} & \textbf{Pipeline} & \textbf{Garmin} & \textbf{Różnica} \\
\hline
Kadencja [spm]   & 162{,}5 & 163{,}3 & $-0{,}5\%$ \\
Czas cyklu [ms]  & 737     & $\approx$735 & $+0{,}3\%$ \\
GCT (śr. L/P) [ms] & 255   & 297     & $-14\%$ \\
\hline
\end{tabular}
\end{table}

Wynik ujawnia wyraźny kontrast, który stanowi empiryczne potwierdzenie analizy z~sekcji~\ref{sec:wplyw-na-metryki}. Kadencja i~czas cyklu zgadzają się z~Garminem niemal idealnie (poniżej 1\% różnicy), natomiast czas kontaktu z~podłożem jest systematycznie zaniżony o~około 14\%. Różnica ta wynika z~natury obu metryk: kadencja jest \textbf{licznikiem} kontaktów --- wystarczy wykryć, że kontakt nastąpił, niezależnie od tego, gdzie dokładnie wypada granica między fazą kontaktu a~lotem. Czas kontaktu jest natomiast \textbf{miarą czasu trwania} fazy STANCE, więc zależy bezpośrednio od położenia tej granicy. Pipeline systematycznie skraca fazę STANCE (zaniżając jednocześnie duty factor: 0{,}345 wobec wartości $\approx$0{,}40 wynikającej z~pomiaru Garmina), co wynika z~przyjętego w~auto-etykietowaniu parametru \texttt{flight\_fraction}~=~0{,}4 oraz z~różnicy definicyjnej (Garmin wyznacza kontakt akcelerometrem, od~uderzenia pięty do~oderwania palców). Sprawdzenie to dostarcza ilościowego dowodu na to, że metryki oparte na zliczaniu zdarzeń są znacznie odporniejsze niż metryki oparte na pomiarze czasu trwania faz.

Nagranie potwierdziło również wniosek ze~Studium~1. Mimo poziomej orientacji i~wysokiej rozdzielczości, kąt stopy w~momencie kontaktu wyniósł $-90{,}5^\circ$ (lewa) i~$-98{,}4^\circ$ (prawa) --- wartości fizycznie niemożliwe, automatycznie oznaczone flagą \texttt{low\_confidence}. Pokazuje to, że pozioma orientacja jest warunkiem koniecznym, ale niewystarczającym wiarygodności wzorca lądowania: kamera musi być również ustawiona prostopadle do toru biegu, czego w~tym nagraniu zabrakło. Współczynniki temporalne, niezależne od perspektywy, pozostały przy tym wiarygodne.

Warto odnotować dwie dodatkowe obserwacje. Po pierwsze, symetria lewo-prawo tego biegu była najlepsza w~całym materiale (maksymalny wskaźnik symetrii Robinsona 9{,}5\%, średni 3{,}1\%), co odpowiada balansowi kroków 49,5/50,5\% raportowanemu przez aplikację Garmin Connect. Po drugie, średnia pewność predykcji klasyfikatora wyniosła 0{,}855 --- najniższa spośród analizowanych biegaczy, choć wciąż powyżej progu ostrzeżenia 0{,}85.

\section{Podsumowanie}
\label{sec:podsumowanie-wrazliwosc}

Trzy studia przypadku potwierdziły trzy postawione hipotezy: współczynniki przestrzenne są wrażliwe na perspektywę kamery (Studium~1), temporalne na częstotliwość klatek (Studium~2), a~najbardziej zawodne przypadki --- wartości wykraczające poza zakres fizjologiczny --- da się wykryć automatycznie (sekcja~\ref{sec:auto-detekcja}). Walidacja na nagraniu własnym z~pomiarem referencyjnym Garmin (sekcja~\ref{sec:walidacja-garmin}) dodatkowo potwierdziła te wnioski niezależnym ground-truthem: kadencja zgodna w~granicach 1\%, czas kontaktu zaniżony o~14\% (potwierdzenie wrażliwości metryk czasu trwania), a~wzorzec lądowania automatycznie oznaczony jako niewiarygodny mimo poprawnej rozdzielczości i~orientacji.

Wkładem tej części pracy jest systematyczna analiza wrażliwości pipeline'u monocular 2D dla \textit{biegu}. Literatura badała podobne ograniczenia dla \textit{chodu} --- Stenum i~wsp.~\cite{stenum2021} opisali zależność długości kroku od pozycji w~kadrze, a~Menychtas i~wsp.~\cite{menychtas2023} błędy dla ruchów poza płaszczyzną obrazu --- ale dla biegu sportowego brakowało takiego zestawienia wraz z~działającym mechanizmem automatycznego ostrzegania użytkownika.
